\documentclass[a4paper,11pt]{scrartcl}

\usepackage[T1]{fontenc}
\usepackage[utf8]{inputenc}
\usepackage[ngerman]{babel}
\usepackage{lmodern}
\usepackage{amsmath}
\usepackage{amssymb}
\usepackage[a4paper,margin=2.5cm]{geometry}
\usepackage{enumitem}
\usepackage{booktabs}
\usepackage[hidelinks]{hyperref}
\usepackage{catppuccin-dark}

\newcounter{aufgabe}
\newcommand{\Aufgabe}[2]{%
  \refstepcounter{aufgabe}%
  \subsection*{Aufgabe~\theaufgabe: #1 \hfill (#2~Punkte)}%
  \label{auf:\theaufgabe}}

\title{Übungsblatt 5: Global Optimization (Globale Optimierung)}
\subtitle{Begleitmaterial zur Vorlesung \glqq Numerical Methods\grqq{} (Prof.\ Dr.-Ing.\ Mark Schutera)}
\author{Shekel's Foxholes, Newton für Minima, Random Restarts,\\ Basin Hopping und Simulated Annealing (Skript S.~47--56)}
\date{\today}

\begin{document}

\maketitle

\noindent\textbf{Gesamt: 50 Punkte.}\quad
Hilfsmittel: Taschenrechner. Runden Sie Zwischenergebnisse, sofern nicht anders angegeben, auf
mindestens vier Nachkommastellen; bei Endergebnissen genügen zwei bis drei signifikante
Nachkommastellen. Die Aufgaben steigen im Schwierigkeitsgrad an: Aufgaben~1--3 prüfen das
Begriffsverständnis, Aufgaben~4--7 sind Rechenaufgaben analog zum Skript, Aufgabe~8 ist eine
Transferaufgabe. Die Gleichungsnummern (35)--(41) beziehen sich auf das Skript.

\medskip
\noindent\textbf{Durchgängiges Beispiel:} Mehrere Aufgaben verwenden die konkrete
1D-Shekel-Funktion (Shekel's Foxholes) des Skripts [S.~48] mit drei Fuchslöchern (foxholes),
ein Spezialfall von Gl.~(35), $f(x) = -\sum_{i=1}^{m} \frac{c_i}{(x-a_i)^2 + r_i}$:
\[
  f(x) \;=\; -\frac{1}{(x-0)^2 + 0.7} \;-\; \frac{1.7}{(x-4)^2 + 0.2} \;-\; \frac{1.1}{(x-7)^2 + 0.4}.
\]
Die Funktion besitzt drei Minima bei $x = 0$, $x = 4$ und $x = 7$; das globale Minimum
liegt bei $x^* = 4$.

\bigskip
\hrule

\section*{Aufgaben}

% ------------------------------------------------------------
\Aufgabe{Lokale versus globale Optimierung}{4}

Das Skript unterscheidet [S.~51] die lokale Optimierung (local optimization), Gl.~(37),
\[
  \mathbf{x}^* = \arg\min_{\mathbf{x} \in \mathcal{N}(\mathbf{x}_0)} f(\mathbf{x}), \tag{37}
\]
von der globalen Optimierung (global optimization), Gl.~(38),
\[
  \mathbf{x}^* = \arg\min_{\mathbf{x} \in \Omega} f(\mathbf{x}). \tag{38}
\]

\begin{enumerate}[label=(\alph*)]
  \item Erklären Sie in eigenen Worten den Unterschied zwischen den beiden Problemstellungen.
        Was bezeichnen $\mathcal{N}(\mathbf{x}_0)$ und $\Omega$? Geben Sie eine konkrete
        mengentheoretische Beschreibung einer Nachbarschaft $\mathcal{N}(\mathbf{x}_0)$ an.
        \hfill (2~P)
  \item Warum kann eine lokale Methode wie Newtons Verfahren das Finden des globalen
        Minimums nicht garantieren? Belegen Sie Ihre Antwort mit dem Verhalten auf der
        Shekel-Funktion für die Startpunkte $x_0 = 3$, $x_0 = -2$ und $x_0 = 8$
        (qualitativ, ohne Rechnung). \hfill (2~P)
\end{enumerate}

% ------------------------------------------------------------
\schreibraum[7]
\Aufgabe{Konvexität als Trennlinie (convexity is the dividing line)}{6}

Eine Funktion $f$ heißt \emph{konvex} (convex), wenn für alle $x, y$ und alle
$\lambda \in [0,1]$ gilt [S.~51]:
\[
  f(\lambda x + (1-\lambda)y) \;\leq\; \lambda f(x) + (1-\lambda) f(y).
\]

\begin{enumerate}[label=(\alph*)]
  \item Warum ist Konvexität die \glqq Trennlinie\grqq{} der Optimierung? Was gilt für
        konvexe Probleme, und welche drei Landschaftsmerkmale machen nicht-konvexe
        Optimierung laut Skript schwierig? \hfill (1~P)
  \item Zeigen Sie mit der Definition (nicht über die zweite Ableitung), dass
        $f(x) = x^2$ konvex ist. \emph{Hinweis:} Bringen Sie die Differenz
        $\lambda f(x) + (1-\lambda)f(y) - f(\lambda x + (1-\lambda)y)$ auf eine
        offensichtlich nichtnegative Form. \hfill (2~P)
  \item Zeigen Sie, dass die Shekel-Funktion (Instanz vom Deckblatt) \emph{nicht} konvex
        ist, indem Sie die Definition mit $x = 4$, $y = 7$ und $\lambda = \tfrac12$
        widerlegen. Werten Sie dazu $f(4)$, $f(7)$ und $f(5.5)$ aus. Formulieren Sie
        zusätzlich ein allgemeines Argument: Warum kann eine Funktion mit zwei strikten
        lokalen Minima niemals konvex sein? \hfill (3~P)
\end{enumerate}

% ------------------------------------------------------------
\schreibraum[9]
\Aufgabe{Basin of Attraction und globale Strategien im Vergleich}{5}

\begin{enumerate}[label=(\alph*)]
  \item Definieren Sie den Begriff \emph{Anziehungsgebiet} (basin of attraction) eines
        lokalen Minimums im Kontext eines lokalen Optimierers. \hfill (1~P)
  \item Vergleichen Sie die drei globalen Strategien \emph{Random Restarts}
        (Algorithmus~3 [S.~51]), \emph{Basin Hopping} (Algorithmus~4 [S.~52]) und
        \emph{Simulated Annealing} [S.~52] hinsichtlich: (i)~Wie wird der nächste
        Startpunkt bzw.\ Kandidat erzeugt? (ii)~Wird Information aus dem bisherigen
        Suchverlauf weiterverwendet oder die Suche zurückgesetzt? (iii)~Nach welchem
        Kriterium werden Kandidaten akzeptiert? \hfill (3~P)
  \item Was bedeutet der Hyperparameter \emph{Patience} bei der Wahl der Restart-Anzahl
        $K$, wenn die Basin-Wahrscheinlichkeit $p$ vorab unbekannt ist? \hfill (1~P)
\end{enumerate}

% ------------------------------------------------------------
\schreibraum[8]
\Aufgabe{Newton für Minima und der Test der zweiten Ableitung}{7}

Um mit Newtons Methode Optima zu finden, wird das Optimierungsproblem als
Nullstellenproblem auf der Ableitung reformuliert: Wir lösen $f'(x) = 0$ und iterieren
nach Gl.~(36)/(40):
\[
  x_{n+1} = x_n - \frac{f'(x_n)}{f''(x_n)}. \tag{40}
\]
Newtons Methode, angewandt auf $f'$, findet dabei Punkte mit
$f'(x) = 0$\footnote{Korrektur zum Originalskript: Auf S.~53 steht versehentlich
\glqq Newton's method finds points where $f''(x) = 0$\grqq. Sachlich richtig findet das
auf $f'$ angewandte Newton-Verfahren Punkte mit $f'(x) = 0$ (Extremstellen); erst der
anschließende Test der zweiten Ableitung klassifiziert diese über $f''$.} ---
das können Minima, Maxima oder Sattelpunkte sein. Gegeben sei
\[
  g(x) = x^3 - 6x^2 + 9x + 1.
\]

\begin{enumerate}[label=(\alph*)]
  \item Bestimmen Sie $g'(x)$ und $g''(x)$, berechnen Sie alle stationären Punkte
        ($g'(x) = 0$) analytisch und klassifizieren Sie sie mit dem Test der zweiten
        Ableitung (second derivative test). Geben Sie auch die Funktionswerte an den
        stationären Punkten an. Nennen Sie alle drei Fälle des Tests. \hfill (3~P)
  \item Führen Sie ab dem Startpunkt $x_0 = 3.5$ zwei Newton-Schritte nach Gl.~(40)
        aus. Wohin konvergiert das Verfahren offenbar, und um welchen Typ von
        stationärem Punkt handelt es sich? \hfill (2~P)
  \item Führen Sie ab dem Startpunkt $x_0 = 0.5$ einen Newton-Schritt aus. Wohin läuft
        das Verfahren hier? Was folgt daraus für die Aussage \glqq Newton auf $f'$
        findet Minima\grqq? \hfill (2~P)
\end{enumerate}

% ------------------------------------------------------------
\schreibraum[10]
\Aufgabe{Erfolgswahrscheinlichkeit von Random Restarts}{7}

Liegt der Startpunkt mit Wahrscheinlichkeit $p$ im Anziehungsgebiet des globalen
Minimums, so gilt für $K$ unabhängige Restarts [S.~51], Gl.~(39):
\[
  P(\text{global gefunden}) = 1 - (1-p)^K. \tag{39}
\]

\begin{enumerate}[label=(\alph*)]
  \item \textbf{Skript-Rechnung [S.~54]:} Eine Funktion hat fünf lokale Minima; das
        Basin des globalen Minimums deckt $15\,\%$ der Suchdomäne ab ($p = 0.15$).
        Berechnen Sie die Erfolgswahrscheinlichkeit für $K = 10$ Restarts. \hfill (2~P)
  \item Wie viele Restarts $K$ sind für eine Erfolgswahrscheinlichkeit von $99\,\%$
        nötig? Leiten Sie zunächst die allgemeine Formel für $K$ aus Gl.~(39) her und
        setzen Sie dann ein. \hfill (2~P)
  \item \textbf{Neue Zahlen:} Nun gelte $p = 0.2$. Berechnen Sie (i) die
        Erfolgswahrscheinlichkeit für $K = 8$ Restarts und (ii) die minimale ganzzahlige
        Restart-Anzahl $K$ für mindestens $95\,\%$ Erfolgswahrscheinlichkeit
        (Probe nicht vergessen). \hfill (2~P)
  \item \textbf{Folgeaufgabe des Skripts [S.~55]:} Berechnen Sie $K$ für $p = 0.05$ und
        $p = 0.01$ bei gleichem $99\,\%$-Ziel. Wie skaliert der Aufwand näherungsweise
        mit der Basin-Größe $p$? \hfill (1~P)
\end{enumerate}

% ------------------------------------------------------------
\schreibraum[10]
\Aufgabe{Simulated Annealing: Akzeptanzwahrscheinlichkeit}{5}

Bei Simulated Annealing [S.~52] wird ein Sprung in eine \emph{schlechtere} Lösung
(Verschlechterung $\Delta f > 0$ des Zielfunktionswerts) mit Wahrscheinlichkeit
\[
  P_{\text{akzeptiert}} = \exp\!\left(-\frac{\Delta f}{T}\right)
\]
akzeptiert, wobei die Temperatur $T$ über die Zeit abnimmt.

\begin{enumerate}[label=(\alph*)]
  \item Werten Sie die Akzeptanzwahrscheinlichkeit aus für
        \begin{enumerate}[label=(\roman*)]
          \item $\Delta f = 0.5$ bei den Temperaturen $T = 2$, $T = 0.5$ und $T = 0.1$,
          \item $T = 1$ bei den Verschlechterungen $\Delta f = 0.1$ und $\Delta f = 1.2$.
        \end{enumerate} \hfill (3~P)
  \item Diskutieren Sie anhand Ihrer Zahlen den Temperatureffekt: Wie verhält sich der
        Algorithmus bei hoher bzw.\ niedriger Temperatur, und warum ist dieses Verhalten
        (Exploration zu Beginn, Konservativität am Ende) für globale Optimierung
        sinnvoll? Wie werden Kandidaten mit $\Delta f \leq 0$ (Verbesserungen)
        behandelt? \hfill (2~P)
\end{enumerate}

% ------------------------------------------------------------
\schreibraum[8]
\Aufgabe{Newton Richtung Minima lenken: der $|f''|$-Trick}{6}

Das Skript schlägt [S.~53] das modifizierte Update vor, Gl.~(41):
\[
  x_{n+1} = x_n - \frac{f'(x_n)}{|f''(x_n)|}. \tag{41}
\]

\begin{enumerate}[label=(\alph*)]
  \item Erklären Sie, welches Problem des Standard-Updates Gl.~(40) der $|f''|$-Trick
        behebt. In welchen Regionen der Loss-Landschaft unterscheiden sich die beiden
        Updates? \hfill (1~P)
  \item Demonstrieren Sie den Unterschied an $f(x) = x^4 - 2x^2$ mit Startpunkt
        $x_0 = 0.3$: Bestimmen Sie $f'(x)$, $f''(x)$ sowie alle stationären Punkte mit
        ihrem Typ. Berechnen Sie dann je einen Schritt nach Gl.~(40) und nach Gl.~(41)
        und vergleichen Sie die Funktionswerte $f(x_0)$ und $f(x_1)$ für beide
        Varianten. Welcher Schritt läuft bergab (downhill), welcher bergauf --- und
        wohin würde jede Variante konvergieren? \hfill (3~P)
  \item Was bedeutet der Trick geometrisch für die Schrittrichtung relativ zum
        Gradienten? Und: \emph{How would you modify the update to find maxima instead?}
        --- Geben Sie ein Update an, das stets in Richtung \emph{wachsender}
        Funktionswerte läuft, und nennen Sie eine äquivalente Alternative. \hfill (2~P)
\end{enumerate}

% ------------------------------------------------------------
\schreibraum[9]
\Aufgabe{Transfer: Newton auf Shekel's Foxholes, Basins und Deep Learning}{10}

Betrachten Sie die Shekel-Instanz vom Deckblatt. Für Teilaufgabe~(b) dürfen Sie die
zweite Ableitung als gegeben verwenden:
\[
  f''(x) = \frac{2\,(0.7 - 3x^2)}{(x^2 + 0.7)^3}
         + \frac{3.4\,(0.2 - 3(x-4)^2)}{((x-4)^2 + 0.2)^3}
         + \frac{2.2\,(0.4 - 3(x-7)^2)}{((x-7)^2 + 0.4)^3}.
\]

\begin{enumerate}[label=(\alph*)]
  \item Leiten Sie die erste Ableitung $f'(x)$ her. \emph{Hinweis:} Differenzieren Sie
        zunächst allgemein einen Summanden $-\dfrac{c}{(x-a)^2 + r}$ und setzen Sie dann
        die drei Parametersätze ein. \hfill (2~P)
  \item Werten Sie $f'(7.2)$ und $f''(7.2)$ aus (alle drei Summanden einzeln angeben,
        dann summieren). \hfill (3~P)
  \item Berechnen Sie den Newton-Schritt $x_1$ nach Gl.~(36)/(40) ab $x_0 = 7.2$. Die
        weitere Iteration liefert $x_2 \approx 7.02$ und $x_3 \approx 7.00$.
        Interpretieren Sie das Ergebnis: Wohin konvergiert das Verfahren, und was sagt
        das über den Startpunkt $x_0 = 7.2$ im Verhältnis zum globalen Minimum
        $x^* = 4$ aus? \hfill (2~P)
  \item Nehmen Sie an, ein lokaler Optimierer konvergiert stets in das Minimum
        desjenigen Basins, in dem er startet; die Basin-Grenzen liegen an den lokalen
        \emph{Maxima} von $f$, näherungsweise bei
        $x \approx 1.7$ und $x \approx 5.7$. Geben Sie die Anziehungsgebiete (basins of
        attraction) der drei Minima $x = 0, 4, 7$ an. Basin Hopping springe nun vom
        lokalen Minimum $x^\circ = 7$ mit $\delta \sim \mathrm{Uniform}(-3, 3)$:
        Mit welcher Wahrscheinlichkeit landet ein einzelner Sprung im Basin des
        globalen Minimums? \hfill (2~P)
  \item Warum kann Basin Hopping hier das globale Minimum schneller finden als Random
        Restarts? Und der Deep-Learning-Bezug [S.~51]: Warum sind mehrere
        Trainingsläufe eines neuronalen Netzes mit verschiedenen Seeds
        \emph{implizite Random Restarts}, und warum ist diese Strategie für große
        Modelle nicht praktikabel? \hfill (1~P)
\end{enumerate}

\schreibraum[13]
\newpage

% ============================================================
\section*{Lösungen --- erst nach Bearbeitung lesen!}

% ------------------------------------------------------------
\subsection*{Lösung zu Aufgabe~\ref{auf:1}}

\textbf{(a)} Die lokale Optimierung, Gl.~(37), sucht das beste $\mathbf{x}$ nur in einer
\emph{Nachbarschaft} $\mathcal{N}(\mathbf{x}_0)$ des Startpunkts $\mathbf{x}_0$ --- sie
findet also dasjenige Minimum, in dessen Einzugsbereich der Startpunkt liegt. Die globale
Optimierung, Gl.~(38), sucht dagegen über die \emph{gesamte Suchdomäne} $\Omega$ und
verlangt das beste aller Minima. $\mathcal{N}(\mathbf{x}_0)$ ist eine Menge von Punkten um
$\mathbf{x}_0$, z.\,B.
$\mathcal{N}(\mathbf{x}_0) = \{\mathbf{x} : \|\mathbf{x} - \mathbf{x}_0\| < \epsilon\}$
für einen Radius $\epsilon$ [Randnotiz S.~51]; $\Omega$ ist die gesamte zulässige
Suchdomäne, die leicht unendlich groß werden kann.

\textbf{(b)} Lokale Methoden konvergieren zu derjenigen Lösung, die in der Nähe liegt ---
sie sehen nur Funktionswerte und Gradienten an den ausgewerteten Punkten und können nicht
wissen, ob weiter entfernt ein besseres Minimum existiert. Auf der Shekel-Funktion:
Von $x_0 = 3$ aus hat man \emph{Glück} und landet im Basin des globalen Minimums
$x^* = 4$; von $x_0 = -2$ konvergiert Newton zum lokalen Minimum $x^\circ = 0$, von
$x_0 = 8$ zum lokalen Minimum $x^\circ \approx 7$ --- das globale Minimum wird in beiden
Fällen verfehlt [S.~50]. Selbst mit Konvergenzgarantien garantiert lokale Optimierung also
nur das Finden eines \emph{nahen} Extremums.

% ------------------------------------------------------------
\subsection*{Lösung zu Aufgabe~\ref{auf:2}}

\textbf{(a)} Für konvexe Probleme ist jedes lokale Minimum zugleich das globale Minimum
--- jede lokale Methode findet daher das globale Optimum, lokale Verfahren genügen.
Die Schwierigkeit entsteht erst bei nicht-konvexer Optimierung: Dort enthält die
Landschaft (i) mehrere lokale Minima, (ii) Sattelpunkte (saddle points) und
(iii) Plateaus [S.~51].

\textbf{(b)} Für $f(x) = x^2$ betrachten wir die Differenz beider Seiten der Definition:
\begin{align*}
  \lambda x^2 + (1-\lambda)y^2 - (\lambda x + (1-\lambda)y)^2
  &= \lambda x^2 + (1-\lambda)y^2 - \lambda^2 x^2 - 2\lambda(1-\lambda)xy - (1-\lambda)^2 y^2\\
  &= \lambda(1-\lambda)x^2 - 2\lambda(1-\lambda)xy + (1-\lambda)\bigl(1 - (1-\lambda)\bigr)y^2\\
  &= \lambda(1-\lambda)\,(x^2 - 2xy + y^2)
   = \lambda(1-\lambda)\,(x-y)^2 \;\geq\; 0,
\end{align*}
denn $\lambda \in [0,1]$ impliziert $\lambda(1-\lambda) \geq 0$ und $(x-y)^2 \geq 0$.
Also gilt $f(\lambda x + (1-\lambda)y) \leq \lambda f(x) + (1-\lambda)f(y)$: $x^2$ ist
konvex.

\textbf{(c)} Auswertung der drei Funktionswerte:
\begin{align*}
  f(4) &= -\frac{1}{16 + 0.7} - \frac{1.7}{0.2} - \frac{1.1}{9 + 0.4}
        = -0.0599 - 8.5 - 0.1170 \approx -8.6769,\\
  f(7) &= -\frac{1}{49 + 0.7} - \frac{1.7}{9 + 0.2} - \frac{1.1}{0.4}
        = -0.0201 - 0.1848 - 2.75 \approx -2.9549,\\
  f(5.5) &= -\frac{1}{30.25 + 0.7} - \frac{1.7}{2.25 + 0.2} - \frac{1.1}{2.25 + 0.4}
        = -0.0323 - 0.6939 - 0.4151 \approx -1.1413.
\end{align*}
Mit $\lambda = \tfrac12$ verlangt Konvexität
\[
  f(5.5) \;\leq\; \tfrac12 f(4) + \tfrac12 f(7) = \tfrac12(-8.6769 - 2.9549) \approx -5.8159.
\]
Tatsächlich ist aber $f(5.5) \approx -1.1413 > -5.8159$ --- Widerspruch, die
Shekel-Funktion ist nicht konvex.

\emph{Allgemeines Argument:} Hätte eine konvexe Funktion zwei strikte lokale Minima
$x_1 \neq x_2$, so müsste die Funktion auf der Verbindungsstrecke unterhalb der Sehne
bleiben. Zwischen zwei strikten Minima muss die Funktion aber zwingend wieder ansteigen
(es existiert ein Punkt dazwischen mit Funktionswert oberhalb beider Minima, ein lokales
Maximum), wodurch die Sehnenbedingung dort verletzt wird. Bei konvexen Funktionen ist die
Menge der Minimalstellen zusammenhängend (konvex), und jedes lokale Minimum ist global ---
zwei getrennte strikte lokale Minima sind unmöglich.

% ------------------------------------------------------------
\subsection*{Lösung zu Aufgabe~\ref{auf:3}}

\textbf{(a)} Das Anziehungsgebiet (basin of attraction) eines lokalen Minimums ist die
Menge aller Startpunkte $x_0$, von denen aus der lokale Optimierer zu genau diesem
Minimum konvergiert. Die Basin-Grenzen liegen dort, wo die Anziehung von einem Minimum
zum nächsten wechselt (bei 1D-Funktionen typischerweise an den lokalen Maxima zwischen
den Tälern).

\textbf{(b)} Vergleich der drei Strategien:

\begin{center}
\begin{tabular}{@{}p{2.6cm}p{3.7cm}p{3.9cm}p{4.0cm}@{}}
\toprule
 & \textbf{Random Restarts} & \textbf{Basin Hopping} & \textbf{Simulated Annealing}\\
\midrule
(i) Kandidaten- erzeugung &
zufälliger neuer Startpunkt $x_0 \sim \mathrm{RandomPoint}(\Omega)$ aus der gesamten Domäne &
zufälliger Sprung $\delta \sim \mathcal{D}$ (z.\,B. Gauß um null oder uniform) vom zuletzt gefundenen lokalen Optimum: $x \leftarrow x^* + \delta$ &
wie Basin Hopping: Sprung vom aktuellen Punkt aus\\[2pt]
(ii) Suchverlauf &
setzt die Suche komplett zurück (\glqq complete reset\grqq); nur $x_{\text{best}}$ wird gemerkt &
nutzt den bisherigen Verlauf: Sprünge starten beim aktuellen lokalen Optimum, lokale Optimierung wird weiter ausgenutzt &
nutzt den Verlauf wie Basin Hopping\\[2pt]
(iii) Akzeptanz &
implizit: bestes $f(x^*)$ gewinnt am Ende &
bestes $f(x^*)$ wird gemerkt; jeder Sprung wird ausgeführt &
probabilistisch: Verbesserungen immer; Verschlechterungen mit $P = \exp(-\Delta f / T)$, $T$ sinkt über die Zeit\\
\bottomrule
\end{tabular}
\end{center}

Kernunterschied laut Skript [S.~52]: Basin Hopping erlaubt das Entkommen aus lokalen
Minima, während weiterhin lokale Optimierung zum Finden besserer Lösungen genutzt wird;
Random Restarts setzt die Suche dagegen vollständig zurück (\glqq fühlt sich im Ansatz
wieder wie Brute-Forcing an\grqq). Simulated Annealing ergänzt Basin Hopping um das
temperaturgesteuerte Akzeptanzkriterium und wird mit sinkender Temperatur zunehmend
konservativ, bis nur noch verbessernde Kandidaten akzeptiert werden.

\textbf{(c)} \emph{Patience} kontrolliert, wie viele Restarts ohne Verbesserung der
besten gefundenen Lösung abgewartet werden, bevor gestoppt wird: Man erhöht $K$ so lange,
bis sich $x_{\text{best}}$ stabilisiert (keine Verbesserung nach mehreren aufeinander
folgenden Restarts) --- eine Heuristik für die Wahl von $K$, wenn $p$ unbekannt ist
[S.~51].

% ------------------------------------------------------------
\subsection*{Lösung zu Aufgabe~\ref{auf:4}}

\textbf{(a)} Ableitungen:
\[
  g'(x) = 3x^2 - 12x + 9 = 3(x-1)(x-3), \qquad g''(x) = 6x - 12.
\]
Stationäre Punkte: $g'(x) = 0 \Leftrightarrow x = 1$ oder $x = 3$.
Test der zweiten Ableitung:
\begin{itemize}[itemsep=1pt]
  \item $x = 1$: $g''(1) = -6 < 0$ $\Rightarrow$ lokales \textbf{Maximum} (Kurve biegt
        nach unten), $g(1) = 1 - 6 + 9 + 1 = 5$.
  \item $x = 3$: $g''(3) = 6 > 0$ $\Rightarrow$ lokales \textbf{Minimum} (Kurve biegt
        nach oben), $g(3) = 27 - 54 + 27 + 1 = 1$.
\end{itemize}
Die drei Fälle des Tests [S.~53]: $f''(x^*) > 0$ lokales Minimum; $f''(x^*) < 0$ lokales
Maximum; $f''(x^*) = 0$ nicht schlüssig (möglicher Sattelpunkt oder Wendepunkt).

\textbf{(b)} Ab $x_0 = 3.5$:
\[
  g'(3.5) = 3(2.5)(0.5) = 3.75, \quad g''(3.5) = 9
  \;\Rightarrow\; x_1 = 3.5 - \frac{3.75}{9} = 3.5 - 0.4167 = 3.0833.
\]
\[
  g'(3.0833) = 3(2.0833)(0.0833) = 0.5208, \quad g''(3.0833) = 6.5
  \;\Rightarrow\; x_2 = 3.0833 - \frac{0.5208}{6.5} = 3.0833 - 0.0801 = 3.0032.
\]
Das Verfahren konvergiert offensichtlich gegen $x = 3$, das lokale \textbf{Minimum}
($g''(3) = 6 > 0$).

\textbf{(c)} Ab $x_0 = 0.5$:
\[
  g'(0.5) = 3(-0.5)(-2.5) = 3.75, \quad g''(0.5) = -9
  \;\Rightarrow\; x_1 = 0.5 - \frac{3.75}{-9} = 0.5 + 0.4167 = 0.9167.
\]
Das Verfahren läuft auf $x = 1$ zu --- das lokale \textbf{Maximum} (ein weiterer Schritt
liefert $x_2 \approx 0.9968$). Folgerung: Newton auf $f'$ findet \emph{alle} Punkte mit
$f'(x) = 0$, gleichgültig ob Minimum, Maximum oder Sattelpunkt --- das Verfahren
unterscheidet den Typ nicht. Erst der Test der zweiten Ableitung klassifiziert das
Ergebnis; wer gezielt Minima will, braucht z.\,B. den $|f''|$-Trick aus Aufgabe~\ref{auf:7}.

% ------------------------------------------------------------
\subsection*{Lösung zu Aufgabe~\ref{auf:5}}

\textbf{(a)} Mit $p = 0.15$ und $K = 10$ [S.~54]:
\[
  P = 1 - (1 - 0.15)^{10} = 1 - 0.85^{10} = 1 - 0.1969 = 0.803.
\]
Etwa $80\,\%$ Chance --- \glqq not too bad, but far from certain\grqq.

\textbf{(b)} Auflösen von $P = 1 - (1-p)^K$ nach $K$:
$(1-p)^K = 1 - P \Rightarrow K \ln(1-p) = \ln(1-P)$, also
\[
  K = \frac{\ln(1-P)}{\ln(1-p)}
    = \frac{\ln(0.01)}{\ln(0.85)}
    = \frac{-4.6052}{-0.1625} \approx 28.3
  \;\Rightarrow\; \textbf{K = 29 Restarts} \text{ (aufrunden).}
\]

\textbf{(c)} Mit $p = 0.2$:
\begin{itemize}[itemsep=1pt]
  \item[(i)] $K = 8$: $P = 1 - 0.8^{8} = 1 - 0.1678 = 0.832$, also etwa $83\,\%$.
  \item[(ii)] $K = \dfrac{\ln(0.05)}{\ln(0.8)} = \dfrac{-2.9957}{-0.2231} \approx 13.43
        \Rightarrow K = 14$.\\
        Probe: $1 - 0.8^{14} = 1 - 0.0440 = 0.956 \geq 0.95$~\checkmark, während
        $K = 13$ nur $1 - 0.8^{13} = 0.945 < 0.95$ liefert.
\end{itemize}

\textbf{(d)} Bei gleichem $99\,\%$-Ziel:
\[
  p = 0.05:\; K = \frac{\ln(0.01)}{\ln(0.95)} = \frac{-4.6052}{-0.0513} \approx 89.8
  \Rightarrow K = 90,
  \qquad
  p = 0.01:\; K = \frac{\ln(0.01)}{\ln(0.99)} = \frac{-4.6052}{-0.0101} \approx 458.2
  \Rightarrow K = 459.
\]
Wegen $\ln(1-p) \approx -p$ für kleine $p$ gilt $K \approx \ln(1-P)^{-1}\cdot$\,konstant$/p$,
d.\,h. der Aufwand skaliert näherungsweise mit $1/p$: Halbiert sich das Basin des globalen
Minimums, verdoppelt sich (ungefähr) die nötige Restart-Anzahl. In der Praxis ist $p$
unbekannt --- dann hilft die Patience-Heuristik aus Aufgabe~\ref{auf:3}(c).

% ------------------------------------------------------------
\subsection*{Lösung zu Aufgabe~\ref{auf:6}}

\textbf{(a)} Auswertung von $P = \exp(-\Delta f / T)$:

\begin{center}
\begin{tabular}{@{}llll@{}}
\toprule
$\Delta f$ & $T$ & $-\Delta f / T$ & $P_{\text{akzeptiert}}$\\
\midrule
$0.5$ & $2$   & $-0.25$ & $e^{-0.25} \approx 0.779$\\
$0.5$ & $0.5$ & $-1$    & $e^{-1} \approx 0.368$\\
$0.5$ & $0.1$ & $-5$    & $e^{-5} \approx 0.0067$\\
\midrule
$0.1$ & $1$   & $-0.1$  & $e^{-0.1} \approx 0.905$\\
$1.2$ & $1$   & $-1.2$  & $e^{-1.2} \approx 0.301$\\
\bottomrule
\end{tabular}
\end{center}

\textbf{(b)} \emph{Temperatureffekt:} Bei hoher Temperatur ($T = 2$) wird dieselbe
Verschlechterung $\Delta f = 0.5$ mit $\approx 78\,\%$ akzeptiert, bei $T = 0.1$ nur noch
mit $\approx 0.7\,\%$ --- der Algorithmus ist anfangs \emph{explorativ} (springt häufig
auch in schlechtere Lösungen und kann so Basins lokaler Minima verlassen) und wird mit
sinkender Temperatur zunehmend \emph{konservativ}, bis er praktisch nur noch verbessernde
Kandidaten akzeptiert und zu einem Optimum konvergiert [S.~52]. Außerdem zeigt die
zweite Tabellenhälfte: Bei fester Temperatur werden \emph{kleine} Verschlechterungen
($\Delta f = 0.1$: $\approx 90\,\%$) viel eher akzeptiert als große
($\Delta f = 1.2$: $\approx 30\,\%$). Verbesserungen ($\Delta f \leq 0$) werden stets
akzeptiert (formal wäre $\exp(-\Delta f/T) \geq 1$). Dieses Schema balanciert Exploration
(global suchen, solange $T$ hoch) und Exploitation (lokal verfeinern, wenn $T$ klein) ---
genau das, was globale Optimierung braucht.

% ------------------------------------------------------------
\subsection*{Lösung zu Aufgabe~\ref{auf:7}}

\textbf{(a)} Das Standard-Update Gl.~(40) sucht Nullstellen von $f'$ --- es konvergiert
daher genauso bereitwillig zu Maxima und Sattelpunkten wie zu Minima
(vgl.\ Aufgabe~\ref{auf:4}(c)). Die beiden Updates unterscheiden sich genau dort, wo
$f''(x_n) < 0$ ist (konkave Regionen, z.\,B. in der Nähe lokaler Maxima): Dort kehrt
Gl.~(41) das Vorzeichen des Schritts um. Wo $f''(x_n) > 0$ gilt, sind beide Updates
identisch.

\textbf{(b)} $f(x) = x^4 - 2x^2$, also $f'(x) = 4x^3 - 4x = 4x(x-1)(x+1)$ und
$f''(x) = 12x^2 - 4$. Stationäre Punkte: $x = 0$ mit $f''(0) = -4 < 0$ (lokales
\textbf{Maximum}, $f(0) = 0$) sowie $x = \pm 1$ mit $f''(\pm 1) = 8 > 0$ (globale
\textbf{Minima}, $f(\pm 1) = -1$).

Am Startpunkt $x_0 = 0.3$:
\[
  f'(0.3) = 4(0.027) - 1.2 = -1.092, \qquad f''(0.3) = 12(0.09) - 4 = -2.92 < 0.
\]
\emph{Standard-Newton, Gl.~(40):}
\[
  x_1 = 0.3 - \frac{-1.092}{-2.92} = 0.3 - 0.3740 = -0.0740.
\]
Funktionswerte: $f(0.3) = 0.0081 - 0.18 = -0.1719$, aber
$f(-0.0740) \approx 0.00003 - 0.01095 = -0.0109$. Der Funktionswert ist
\emph{gestiegen}: Der Schritt läuft bergauf, direkt auf das lokale Maximum $x = 0$ zu ---
dorthin würde Standard-Newton konvergieren.

\emph{Mit $|f''|$, Gl.~(41):}
\[
  x_1 = 0.3 - \frac{-1.092}{|-2.92|} = 0.3 + 0.3740 = 0.6740.
\]
Funktionswert: $f(0.6740) \approx 0.2063 - 0.9085 = -0.7022 < -0.1719$. Der Schritt
läuft bergab in Richtung des Minimums $x = 1$ --- dorthin würde die modifizierte
Iteration konvergieren.

\textbf{(c)} \emph{Geometrie:} Da $|f''(x_n)| > 0$, hat der Schritt
$-f'(x_n)/|f''(x_n)|$ stets das \emph{entgegengesetzte} Vorzeichen zu $f'(x_n)$ --- er
zeigt also immer in Richtung des negativen Gradienten (\glqq downhill so to speak\grqq),
während die Krümmung nur noch die Schrittweite skaliert. So bewegen wir uns garantiert in
Richtung abnehmender Funktionswerte $f(x)$ [S.~53]. Für die Tangentenkonstruktion an
$f'$ bedeutet das: In konkaven Regionen wird die Steigung der Tangente am Graphen von
$f'$ effektiv am Vorzeichen gespiegelt, sodass der Schnittpunkt mit der $x$-Achse auf der
bergab gelegenen Seite gesucht wird.

\emph{Maxima statt Minima:} Vorzeichen des Schritts umdrehen,
\[
  x_{n+1} = x_n + \frac{f'(x_n)}{|f''(x_n)|},
\]
dann läuft die Iteration stets in Richtung des \emph{positiven} Gradienten (bergauf,
hill climbing). Äquivalente Alternative: Gl.~(41) unverändert auf $-f$ anwenden, denn
Maxima von $f$ sind Minima von $-f$.

% ------------------------------------------------------------
\subsection*{Lösung zu Aufgabe~\ref{auf:8}}

\textbf{(a)} Für einen allgemeinen Summanden gilt mit Ketten- und Potenzregel:
\[
  \frac{d}{dx}\left[-\frac{c}{(x-a)^2 + r}\right]
  = -c \cdot \frac{d}{dx}\bigl((x-a)^2 + r\bigr)^{-1}
  = \frac{2c\,(x-a)}{\bigl((x-a)^2 + r\bigr)^2}.
\]
Einsetzen der drei Parametersätze $(a, c, r) = (0, 1, 0.7),\ (4, 1.7, 0.2),\ (7, 1.1, 0.4)$:
\[
  f'(x) = \frac{2x}{(x^2 + 0.7)^2}
        + \frac{3.4\,(x-4)}{((x-4)^2 + 0.2)^2}
        + \frac{2.2\,(x-7)}{((x-7)^2 + 0.4)^2}.
\]

\textbf{(b)} Erste Ableitung bei $x = 7.2$ (Summanden einzeln):
\[
  \frac{2 \cdot 7.2}{(51.84 + 0.7)^2} = \frac{14.4}{2760.45} \approx 0.0052, \qquad
  \frac{3.4 \cdot 3.2}{(10.24 + 0.2)^2} = \frac{10.88}{108.99} \approx 0.0998, \qquad
  \frac{2.2 \cdot 0.2}{(0.04 + 0.4)^2} = \frac{0.44}{0.1936} \approx 2.2727,
\]
\[
  f'(7.2) \approx 0.0052 + 0.0998 + 2.2727 = 2.3777 \approx 2.38.
\]
(Das Skript [S.~49] schreibt $0.005 + 0.100 + 2.27 = 2.38$; die gerundeten Summanden
ergeben arithmetisch $2.375$, die Rundung auf $2.38$ ist zulässig und konsistent mit dem
exakten Wert $2.3778$.)

Zweite Ableitung bei $x = 7.2$ (Summanden einzeln):
\[
  \frac{2\,(0.7 - 155.52)}{52.54^3} = \frac{-309.64}{145034} \approx -0.0021, \qquad
  \frac{3.4\,(0.2 - 30.72)}{10.44^3} = \frac{-103.768}{1137.89} \approx -0.0912, \qquad
  \frac{2.2\,(0.4 - 0.12)}{0.44^3} = \frac{0.616}{0.08518} \approx 7.2314,
\]
\[
  f''(7.2) \approx -0.0021 - 0.0912 + 7.2314 = 7.1381 \approx 7.14.
\]

\textbf{(c)} Newton-Schritt nach Gl.~(36)/(40):
\[
  x_1 = 7.2 - \frac{f'(7.2)}{f''(7.2)} = 7.2 - \frac{2.3777}{7.1381}
      = 7.2 - 0.3331 \approx 6.87.
\]
Mit $x_2 \approx 7.02$ und $x_3 \approx 7.00$ konvergiert das Verfahren zum
\emph{lokalen} Minimum $x^\circ \approx 7$ (dort ist $f'' > 0$, der dritte Summand
dominiert). Interpretation: Der Startpunkt $x_0 = 7.2$ liegt im Anziehungsgebiet des
Fuchslochs bei $x = 7$ --- \glqq Here, we were unlucky, and the method is attracted to
the nearby foxhole at $x^\circ \approx 7$\grqq{} [S.~50]. Das globale Minimum $x^* = 4$
wird verfehlt, obwohl es deutlich tiefer liegt ($f(4) \approx -8.68$ gegenüber
$f(7) \approx -2.95$, vgl.\ Aufgabe~\ref{auf:2}(c)): Lokale Optimierung garantiert nur
das Finden eines nahen Extremums.

\textbf{(d)} Da der Optimierer stets in das Minimum seines Start-Basins konvergiert und
die lokalen Maxima die Basin-Grenzen bilden, gilt näherungsweise:
\[
  B(0) = (-\infty,\ 1.7), \qquad B(4) = (1.7,\ 5.7), \qquad B(7) = (5.7,\ \infty).
\]
Das Basin des globalen Minimums $x^* = 4$ ist das mittlere Intervall.
Ein Sprung von $x^\circ = 7$ mit $\delta \sim \mathrm{Uniform}(-3, 3)$ landet
gleichverteilt in $[4, 10]$ (Länge $6$). Der Überlapp mit $B(4) = (1.7, 5.7)$ ist
$[4,\ 5.7)$ mit Länge $1.7$, also
\[
  P = \frac{5.7 - 4}{6} = \frac{1.7}{6} \approx 0.283 \approx 28\,\%.
\]

\textbf{(e)} \emph{Basin Hopping vs.\ Random Restarts:} Basin Hopping startet seine
Sprünge vom bereits gefundenen lokalen Optimum und nutzt damit die Struktur der
Landschaft: Von $x^\circ = 7$ aus erreicht schon ein einzelner Sprung mit
$\approx 28\,\%$ Wahrscheinlichkeit das Basin des globalen Minimums --- die Foxholes
liegen nahe beieinander, die Sprungverteilung muss nur die Distanz zum Nachbar-Basin
überbrücken. Random Restarts sampelt dagegen gleichverteilt aus der gesamten (potenziell
sehr großen) Domäne $\Omega$; ist $\Omega$ groß, wird die Trefferwahrscheinlichkeit $p$
pro Restart klein, und jede Wiederholung verwirft zudem alle bisherige Information
(vollständiger Reset). \emph{Deep-Learning-Bezug [S.~51]:} Jeder Trainingslauf eines
neuronalen Netzes mit anderem Seed startet von einer anderen zufälligen Initialisierung
der Gewichte --- das ist genau ein zufälliger Startpunkt $x_0 \in \Omega$ mit
anschließender lokaler Optimierung (SGD/Adam), also implizite Random Restarts; am Ende
behält man das Modell mit dem besten Loss. Für große Modelle ist das nicht praktikabel,
weil schon ein einzelner Trainingslauf enorm teuer ist --- man verlässt sich stattdessen
auf die Hoffnungen aus dem Skript [S.~53]: Die meisten lokalen Minima sind etwa gleich
gut, Überparametrisierung macht schlechte Minima selten, Mini-Batch-Rauschen hilft beim
Entkommen aus scharfen Minima, adaptive Lernraten helfen auf Plateaus.

\end{document}
